\chapter{Discussion}
\label{ch:discussion}

This chapter answers the three research questions of \S\ref{ch:introduction} and
then draws the forecasting and system threads together. The forecasting answer
(RQ1) rests on results reported in Chapter~\ref{ch:results-models}; the
engineering and maturity answers (RQ2, RQ3) synthesise the system implementation
and its evaluation.

% -----------------------------------------------------------------------------
\section{RQ1 --- Comparing the forecasting models}

The head-to-head comparison of LEAR, the DNN, and the Transformer across DK1 and
DE is reported in Chapter~\ref{ch:results-models}, on a single shared test window
with paired Diebold--Mariano and Giacomini--White tests, and against both a
weekly-naive baseline and the rankings of Lago et al.~\cite{lago2021} and
Llorente and Portela~\cite{llorente2025transformer}. Because the source papers'
test periods predate the present price regime, the comparison is read through
rMAE and model rankings rather than absolute errors.
% NUMBERS PENDING THE SIX --track RUNS (filled in once Chapter ch:results-models is populated).
\textbf{[pending: headline ranking across the three models per market, the rMAE
summary, and whether the GW tests find the differences significant; one paragraph
on what replicated from the papers and what did not.]}

% -----------------------------------------------------------------------------
\section{RQ2 --- Engineering trade-offs of productionising an EPF benchmark}

Embedding the benchmark models in a scheduled, containerised pipeline surfaced a
series of trade-offs, each resolved in favour of the single-host MVP scope. The
most consequential is the choice of a full refit from scratch on every daily run
over warm-start or incremental learning: because the daily data increment is
small relative to the calibration windows, a clean refit is cheap and avoids
accumulated approximation error, and --- as RQ3 and the drift discussion show ---
it also makes the model fresh enough that drift-triggered retraining becomes
redundant. Co-locating training and prediction in a single Prefect flow trades a
small amount of recomputation for the guarantee that the served model always
matches the forecast it produced, removing a class of coordination failure. Using
ClickHouse as the single store for time-series, predictions, metrics, and model
artifacts makes day-after reconciliation one query rather than a cross-system
join, at the cost of a base64-blob artifact representation that will not scale to
large neural weights. And serving the paper-strict feature set in production ---
rather than a separate, richer operational feature set --- keeps the deployed
models bit-faithful to the benchmark, at the cost of forgoing features that might
help operationally but would break the replication. The recurring pattern is that
the constraints of a faithful benchmark and the constraints of a small production
platform mostly align: both reward simplicity, determinism, and a single source
of truth.

% -----------------------------------------------------------------------------
\section{RQ3 --- Critical practices and the limits of a single-host platform}

Of Zarour's nine practices, eight are realised end-to-end and one --- the
feedback loop --- is minimally realised (Table~\ref{tab:maturity}). The
quantitative evaluation locates the binding constraint precisely: the ML Test
Score of 2.5 is bound by the data-testing section, not by infrastructure or
monitoring, which both score higher. In other words, for this platform the
practices that proved easiest to realise to a high level were the
DevOps-inherited ones --- CI/CD automation, versioning, containerised
reproducibility, and monitoring --- because they transfer directly from software
engineering, whereas the ML-specific data-quality discipline (feature-benefit,
feature-cost, and slice analyses) is where the platform is weakest and where
further investment would most raise its maturity. The single-host scope reaches
its limit at three identifiable points: the artifact store, which suits LEAR's
hundred-kilobyte blobs but not multi-megabyte neural weights; the absence of a
model-quality gate before serving; and the coverage gap, with forecasting and
drift scoped to two of the seven ingested bidding zones. None of these is a flaw
in the design so much as a deliberate boundary of the MVP, each with a named
adoption trigger (Chapter~\ref{ch:system-implementation}).

% -----------------------------------------------------------------------------
\section{Tying the threads together}

The two threads are mutually reinforcing rather than independent. Several
modelling decisions are only defensible because of the system around them:
downsampling both markets to an hourly grid is what makes the cross-market and
cross-paper comparison valid; and the cheap, scheduled daily refit is what
licenses treating drift detection as monitoring rather than as a retraining
trigger, following the cost-aware argument of Mahadevan and
Mathioudakis~\cite{mahadevan2024costaware}. Conversely, the system's central
claim --- that any forecast can be reconstructed from its logged provenance ---
is what makes the forecasting replication itself trustworthy: the worked re-run
that reproduced a logged LEAR forecast to within $7\times10^{-6}$\,EUR/MWh
(\S\ref{sec:res-repro}) is evidence about the modelling thread as much as the
system thread. The forecasting payload exercises and validates the platform, and
the platform makes the forecasting reproducible; neither half of the thesis would
stand on its own.

% -----------------------------------------------------------------------------
\section{Limitations}

The limitations span both threads and are collected here. On the modelling side,
the test window is 31 weeks rather than the one-to-two years the source papers
recommend, a consequence of the data history; the DNN hyperparameters are tuned
on an architecture-only search space, with Lago's ten binary feature-selector
hyperparameters left as fixed config decisions because the preprocessor exposes
no toggles for them; the Transformer reuses the paper's search-range defaults
rather than per-market optima; and the work forecasts point prices, not
probabilistic distributions. On the system side, the evaluation is a single-case
study with no controlled comparison, the ML Test Score is self-administered (its
per-item evidence in Appendix~\ref{app:mlts-worksheet} is the mitigation),
forecasting and drift cover only DK1 and DE, the ClickHouse blob artifact store
will not scale to large neural weights, and the drift detector over-flags. These
are revisited as concrete future work in Chapter~\ref{ch:conclusion}, and the
modelling-specific items are detailed alongside the results in
Chapter~\ref{ch:results-models}.
